% New English translation from checked cote 25, site 27 (PDF page 28).
% AI-translated by GPT-6 Astra (OpenAI, ChatGPT Pro mode).
\scanpage{25/27}
\mspage{\hand{23}}

\var{\deleted{Second thoughts on No.4. Results 4.6. b) and c) are entirely
contained in 4.7. b) and c); I had thought such a repetition justifiable
from the standpoint of proof, that of 4.7. seeming to me to use
markedly more delicate results (in the form of 15.2.2.). Now, as long as
we only want to prove in 4.7. b) of c) that U is open, (excluding
the fact $\underline{G}$ U flat \var{\xout{xxx}}{}{typewriter strike-out}{ega25_25_027-v0003} rel to $\overset{\hand{\vee}}{P}$), we have no need at all of 15.2.2.; the
usual method (used in par.12) by the valuative criterion gives the
desired result quite simply. I therefore suggest:a\var{\xout{\unsure{x}}}{}{typewriter strike-out}{ega25_25_027-v0004}) removing
parts b) and c) of 4.6.; b) stating 4.7. b) and c) with only the
conclusion U resp. E open, using the valuative criterion (of course,
4.7. a) is removed from the statement), c) adding another proposition\var{\inserted[typed]{4.8.}}{4.8.}{insertion}{ega25_25_027-v0005}, or a
corollary, saying that under the conditions of 4.7. $\underline{G}$ U is flat rel to $\overset{\hand{\vee}}{P}$,
which is proved with the aid of 15.2.2.}}{}{deletion}{ega25_25_027-v0001}

\var{\over{\unsure{5|6}}{7}}{7}{writing over}{ega25_25_027-v0006}. \uline{Application to the construction of hyperplane sections or multisections
of specified type}.

\var{\over{\unsure{5|6}}{7}}{7}{writing over}{ega25_25_027-v0008}.1. Observe that if S=Spec(k), where k is an \uline{infinite} field, then every nonempty open subset
of $\overset{\hand{\vee}}{P}$ contains a point rational over k; thus with the notation of
\var{\xout{\unsure{xxx}}}{}{typewriter strike-out}{ega25_25_027-v0009}4.1., if E (defined in terms of a constructible property P) contains
the generic point \hand{$\eta$}, it contains \var{\xout{at least}}{}{typewriter strike-out}{ega25_25_027-v0010} a point rational over k, so
there exists a hyperplane section of X (defined over k) having property P.
On the other hand, S being arbitrary again, it is immediate that for every
s\hand{$\in$}S and every point\var{\inserted{$\xi_0$}}{ $\xi_0$}{insertion}{ega25_25_027-v0011} of the fiber $\overset{\hand{\vee}}{P}_s$ rational over k(s), there exists a section
\hand{$\xi$} of $\overset{\hand{\vee}}{P}$ over an open neighborhood U of s that passes through \hand{$\xi$}$_0$. If E is then
again \var{\xout{xxx}}{}{typewriter strike-out}{ega25_25_027-v0012} defined\var{\xout{x}}{}{typewriter strike-out}{ega25_25_027-v0013} as in 4.1. in terms of a constructible property
P, and if we know (e.g. thanks \var{\deleted{to 5.\unsure{5}.}}{}{deletion}{ega25_25_027-v0014}\var{\inserted{to No.5}}{to No.5}{insertion}{ega25_25_027-v0015}) that E is open, then
if \hand{$\xi$}$_0$\hand{$\in$} E, \var{\xout{we can take M}}{}{typewriter strike-out}{ega25_25_027-v0016} the section \hand{$\xi$} is a section of E over \var{\xout{\unsure{Sx}}}{}{typewriter strike-out}{ega25_25_027-v0017}
U, at least if we shrink U sufficiently. Thus we shall have constructed a
hyperplane section\var{\inserted[typed]{$Y_{\hand{\xi}}$}}{ $Y_{\hand{\xi}}$}{insertion}{ega25_25_027-v0018} of X, over an open neighborhood U of s, such that
\tnote{25-027-1}{Remords is rendered second thoughts. In the cancelled paragraph the source has b) de c), rendered b) of c), and writes the underlined G followed by U without a restriction sign; both features are retained. The cancelled stray letter \unsure{r|m} in entièrement does not change English entirely; its doubt is retained here without an English variant. Vider is rendered removing in the instructions to the editor.}

\var{\margin{27}}{\placedmargin{27}}{marginal note placed}{ega25_25_027-v0019}
